\chapter{Solutions}

\runinhead{Dans un triangle
(1)}\phantomsection\addcontentsline{toc}{subsubsection}{Dans un triangle
(1)}\label{III-answer-etb1}
La densité marginale de \(X\) est donnée par
\(f_X(x) = \int f_{X,Y}(x,y) dy = 1_{]0,1[}(x)\) et pour
\(x \in ]0,1[\), \[f_{Y|X=x} (y) = \frac{1}{x} 1_{]0,x[}(y)\] Ainsi
\(X\) est uniformément distribué sur \(]0,1[\), et la loi de \(Y\)
sachant \(X =x\) est uniforme sur \(]0,x[\) pour \((0 < x < 1)\).

\runinhead{Dans un triangle
(2)}\phantomsection\addcontentsline{toc}{subsubsection}{Dans un triangle
(2)}\label{III-answer-etb2}
Pour un tel \(x\), l'espérance conditionnelle \(\mathbb{E}(Y|X=x)\) vaut
ainsi \(x/2\) et nous obtenons \(\mathbb{E}(Y|X) = \frac{X}{2}\).

\runinhead{Auto-conditionnement}\phantomsection\addcontentsline{toc}{subsubsection}{Auto-conditionnement}\label{III-answer-autocond}
On a \(\psi(y) = \mathbb{E}(Y|Y=y) = y\) et donc
\(\mathbb{E}(Y|Y) = \psi(Y) = Y\) p.s.

\runinhead{Espérance conditionnelle d'un produit de
variables}\phantomsection\addcontentsline{toc}{subsubsection}{Espérance
conditionnelle d'un produit de variables}\label{III-answer-prod}
On a \(\mathbb{E}(XY|X = x) = x\mathbb{E}(Y|X=x)\), d'où
\(\mathbb{E}(XY|X) = X\mathbb{E}(Y|X)\) p.s.

\runinhead{Une deuxième
approche}\phantomsection\addcontentsline{toc}{subsubsection}{Une
deuxième approche}\label{III-answer-detail}
On remarque qu'on peut réécrire l'erreur comme
\[\ell(a,b)= \mathbb{E}(Y^2) -2a\mathbb{E}(XY) -2b \mathbb{E}(Y) +a^2\mathbb{E}(X^2) +2ab\mathbb{E}(X) +b^2.\]
Cette fonction est convexe et par conséquent admet un minimum global
atteint au point \((a^*, b^*)\in\mathbb{R}^2\) qui annule les dérivées
dérivées partielles de \(\ell\) en \(a\) et \(b\). Ainsi, on a
\[\frac{\partial \ell(a^*,b^*)}{\partial b} = -2\mathbb{E}(Y^2) + 2a^* \mathbb{E}(X) + 2b^*=0\]
d'où \(b^* = \mathbb{E}(Y) - a^* \mathbb{E}(X)\) Et, \[\begin{aligned}
\frac{\partial \ell(a^*,b^*)}{\partial a} & = -2\mathbb{E}(XY) + 2a \mathbb{E}(X^2) + 2b \mathbb{E}(X)\\
                                   & = -2\mathbb{E}(XY) + 2a \mathbb{E}(X^2) + 2\mathbb{E}(X)\mathbb{E}(Y) - 2 a^* \mathbb{E}(X^2)\\
                                   & = -2\text{Cov}(X,Y) + a^*\mathbb{V}(X)=0
\end{aligned}\]

d'où
\(a^* = \frac{\text{Cov}(X,Y)}{\mathbb{V}(X)} = \rho(X,Y)\frac{\sigma_Y}{\sigma_X}.\)

\runinhead{Variance
totale}\phantomsection\addcontentsline{toc}{subsubsection}{Variance
totale}\label{III-answer-vartot}

\begin{align*}
\mathbb{V}(Y) =& \mathbb{E}((Y-\mathbb{E}(Y))^2) = \mathbb{E}(\mathbb{E}((Y-\mathbb{E}(Y))^2|X)) \text{   par la formule de l'espérance totale}\\
=& \mathbb{E}(\mathbb{E}((Y-\mathbb{E}(Y|X)+\mathbb{E}(Y|X)-\mathbb{E}(Y))^2|X)) \\
=& \mathbb{E}(\mathbb{E}((Y-\mathbb{E}(Y|X))^2|X)) + \mathbb{E}(\mathbb{E}((\mathbb{E}(Y|X)-\mathbb{E}(Y))^2|X)) \\
 &+ 2\mathbb{E}(\mathbb{E}((Y-\mathbb{E}(Y|X))(\mathbb{E}(Y|X)-\mathbb{E}(Y))|X))\\
=& \mathbb{E}(\mathbb{V}(Y|X)) + \mathbb{E}((\mathbb{E}(Y|X)-\mathbb{E}(Y))^2) + 2\mathbb{E}((\mathbb{E}(Y|X)-\mathbb{E}(Y))\mathbb{E}((Y-\mathbb{E}(Y|X))|X))\\
=& \mathbb{E}(\mathbb{V}(Y|X)) + \mathbb{V}(\mathbb{E}(Y|X)) \text{   car    } \mathbb{E}((Y-\mathbb{E}(Y|X))|X) = 0 
\end{align*}

\section{Couple de variables}\label{III-couple-de-variables-1}

\runinhead{Question
1}\phantomsection\addcontentsline{toc}{subsubsection}{Question
1}\label{III-answer-ex1-1}
On voit d'abord que la densité du couple \((X,Y)\) vaut~:
\[f_{X,Y}(x,y) = f_{X|Y=y}(x)f_Y(y) = xe^{-xy}1_{\mathbb{R}_+}(x) 1_{[1,+\infty[} (y)\]
Soit \(h\) une fonction continue bornée sur \(\mathbb{R}^2_+\). Le
changement de variable \((x, y) \mapsto (t = xy , y)\) de jacobien
\(y\), donne alors que
\[\mathbb{E}(h(T, Y )) = \mathbb{E}(h(XY, Y )) = \int_1^{+\infty}\int_0^{+\infty} h(t,y) e^{-t} \frac{t}{y^2} dt dy\]
et donc la densité du couple \((T, Y)\) vaut
\[f_{T,Y} = e^{-t} \frac{t}{y^2}1_{[1,+\infty[} (y)1_{\mathbb{R}_+}(t)\]
Elle s'écrit comme produit d'une fonction de \(t\) et d'une fonction de
\(y\). On en déduit que \(T\) et \(Y\) sont indépendantes et que \(T\) a
pour densité \(te^{-t} 1_{\mathbb{R}_+}(t)\).

\runinhead{Question
2}\phantomsection\addcontentsline{toc}{subsubsection}{Question
2}\label{III-answer-ex1-2}
La loi marginale de \(X\) a pour densité
\(f_X (x) = \int_1^{+\infty} x e^{-xy} dy = e^{-x}\). Ainsi \(X\) suit
une loi exponentielle de paramètre 1 et la loi conditionnelle de \(Y\)
sachant \(X = x\) admet la densité~:
\[f_{Y|X=x}(y) = \frac{f_{X,Y}(x,y)}{f_X(x)} = x e^{-x(y-1)} 1_{[1,+\infty[} (y)\]
pour x \textgreater{} 0.

\runinhead{Question
3}\phantomsection\addcontentsline{toc}{subsubsection}{Question
3}\label{III-answer-ex1-3}
On en déduit que
\[\mathbb{E}(Y|X=x ) = \int_1^{+\infty} y x e^{-x(y-1)} dy = \frac{x+1}{x}1_{\mathbb{R}_+}(x)\]
par intégration par parties. Ainsi \(\mathbb{E}(Y|X) = \frac{X+1}{X}\).

\section{Mélanges de lois}\label{III-muxe9langes-de-lois-1}

\runinhead{Question
1}\phantomsection\addcontentsline{toc}{subsubsection}{Question
1}\label{III-answer-melloi1}
Soit \(B\) un borélien. Par indépendance de \(K\) avec \(X_i\), on a
\[\mathbb{P}(X \in B \mid K = i) = \mathbb{P}(X_i \in B \mid K = i) = \mathbb{P}(X_i \in B).\]
La loi de \(X\) sachant \(\{K = i\}\) est donc la même que celle de
\(X_i\), d'où
\[f_{X\mid K = i} : x\in\mathbb{R}\mapsto f_i(x) = \dfrac{1}{\sqrt{2\pi}\sigma_i}\,\exp\left\{- \dfrac{(x-m_i)^2}{2\sigma_i^2} \right\}.\]

\runinhead{Question
2}\phantomsection\addcontentsline{toc}{subsubsection}{Question
2}\label{III-answer-melloi2}
Soit \(B\) un borélien. D'après la formule des probabilités totales et
la question précédente, on a
\[\mathbb{P}(X\in B) = \sum_{i = 1}^n p_i\,\mathbb{P}(X \in B \mid K = i) = \sum_{i = 1}^n p_i\,\mathbb{P}(X_i \in B).\]
La variable aléatoire \(X\) admet donc une densité, qui vaut
\[f_X : x\in\mathbb{R}\mapsto \sum_{i = 1}^n p_i\,f_i(x).\]

\runinhead{Question
3}\phantomsection\addcontentsline{toc}{subsubsection}{Question
3}\label{III-answer-melloi3}
D'après la question précédente, \(X\) a pour espérance \begin{align*}
\mathbb{E}(X) &= \int_\mathbb{R}x\,f_X(x)\,dx = \int_\mathbb{R}x \sum_{i = 1}^n p_i\, f_i(x)\,dx =  \sum_{i = 1}^n p_i \int_\mathbb{R}x\,f_i(x)\,dx\\
& = \sum_{i = 1}^n p_i\,m_i.
\end{align*} Quant à la variance de \(X\), en utilisant l'égalité
\(\sum_{i= 1}^n p_i = 1\), elle vaut \begin{align*}
\mathbb{V}(X) &= \mathbb{E}\left(X^2\right) - \mathbb{E}(X)^2 = \int_\mathbb{R}x^2\,f_X(x)\,dx - \left(\sum_{i = 1}^n p_i\,m_i\right)^2\\
&= \sum_{i = 1}^n p_i (\sigma_i^2 + m_i^2) - \sum_{i = 1}^n p_i\left(\sum_{j = 1}^n p_j\,m_j\right)^2\\
&= \sum_{i = 1}^n p_i\,\sigma_i^2 + \sum_{i = 1}^n p_i \left(m_i - \sum_{j=1}^n p_j\,m_j \right)^2.
\end{align*} On retrouve bien la forme désirée, avec la dispersion des
espérances
\[\bar{\sigma}^2 := \sum_{i = 1}^n p_i \left(m_i - \sum_{j=1}^n p_j\,m_j \right)^2.\]

\runinhead{Question
4}\phantomsection\addcontentsline{toc}{subsubsection}{Question
4}\label{III-answer-melloi4}
Si l'on souhaite approcher la loi de \(X\) avec une unique Gaussienne,
et non un mélange, les questions précédentes suggèrent de prendre celle
d'espérance \(\sum_{i =1}^n p_i\,m_i\) et de variance
\(\sum_{i = 1}^n p_i\,\sigma^2_i + \bar{\sigma}^2\). Voir figure
ci-dessous.

\begin{figure}
\centering
\includegraphics{PdfMelGauss.tex.pdf}
\caption{Illustration}
\end{figure}

\section{Lois conjuguées}\label{III-answer-loiconj}

On considère dans tout cet exercice \(B_1\) et \(B_2\) des Boréliens.

\runinhead{Question
1}\phantomsection\addcontentsline{toc}{subsubsection}{Question
1}\label{III-answer-loiconj-expexp}
D'après les hypothèses on a \begin{align*}
\mathbb{P}_{X,Y}(B_1\times B_2) &= \int_{B_2} \left(\int_{B_1} d\mathbb{P}_{X\mid Y = y}(x) \right) d\mathbb{P}_Y(y) \hspace{1em}\text{par Fubini conditionnel,}\\
&= \int_{B_2} \left(\int_{B_1} y\,e^{-y x}\,1_{\mathbb{R}_+^\ast}(x)\,dx \right) \lambda\,e^{-\lambda y}\,1_{\mathbb{R}_+^\ast}(y)\,dy\\
&= \int_{B_1} \int_{B_2} \lambda\, y\,e^{-(x+\lambda)\,y}1_{\mathbb{R}_+^\ast}(x)\,1_{\mathbb{R}_+^\ast}(y) \,dy\, dx \hspace{1em}\text{par Fubini.}
\end{align*} Le vecteur aléatoire \((X,Y)\) possède donc une densité
jointe
\[f_{X,Y} : (x,y) \in\mathbb{R}^2 \mapsto \lambda\, y\,e^{-(x+\lambda)\,y}1_{\mathbb{R}_+^\ast}(x)\,1_{\mathbb{R}_+^\ast}(y).\]
La variable aléatoire \(X\) a donc aussi une densité~: pour tout
\(x\in\mathbb{R}\) \begin{align*}
f_X(x) &= \int_{\mathbb{R}} f_{X,Y}(x,y)\,dy = \int_\mathbb{R}\lambda\, y\,e^{-(x+\lambda)\,y}1_{\mathbb{R}_+^\ast}(x)\,1_{\mathbb{R}_+^\ast}(y)\,dy\\
&= \left|\begin{array}{ll}\displaystyle \dfrac{\lambda}{x+\lambda} \int_{0}^{+\infty} y\,(x+\lambda)\,e^{-(x+\lambda)\,y}\,dy & \text{si } x>0,\\[1em] 0 & \text{sinon.} \end{array}\right.
\end{align*} On reconnaît dans cette dernière intégrale la formule de
l'espérance d'une loi Exponentielle de paramètre \(x+\lambda\), et on en
déduit que pour tout \(x\in\mathbb{R}\)
\[f_X(x) = \dfrac{\lambda}{(x+\lambda)^2}\,1_{\mathbb{R}_+^\ast}(x).\]
Pour tout \(x\in\mathbb{R}_+^\ast\) la variable \(Y\) sachant
\(\{X=x\}\) admet donc aussi une densité, que l'on explicite avec la
formule de Bayes~: pour tout \(y\in\mathbb{R}\) \begin{align*}
f_{Y\mid X=x}(y) &= \dfrac{f_{X,Y}(x,y)}{f_X(x)} = \dfrac{\lambda\, y\,e^{-(x+\lambda)\,y}\,1_{\mathbb{R}_+^\ast}(y)}{\dfrac{\lambda}{(x+\lambda)^2}}\\
&= (x+\lambda)^2\,y\,e^{-(x+\lambda)\,y}\,1_{\mathbb{R}_+\ast}(y).
\end{align*} Comme \(\Gamma(2) = 1\), on reconnaît ici la densité d'une
loi Gamma d'indice \(2\) et de paramètre d'échelle \(x+\lambda\).

\runinhead{Question
2}\phantomsection\addcontentsline{toc}{subsubsection}{Question
2}\label{III-answer-loiconj-gampoi}
D'après les hypothèses, en procédant comme précédemment, on a
\begin{align*}
\mathbb{P}_{X,Y}(B_1\times B_2) &= \int_{B_2} \left(\int_{B_1} d\mathbb{P}_{X\mid Y = y}(x) \right) d\mathbb{P}_Y(y)\\
&= \int_{B_2} \left(\sum_{x\in B_1} \dfrac{y^{x}}{x!}\,e^{-y}\,1_{\mathbb{N}}(x) \right) \dfrac{\theta^{\alpha}}{\Gamma(\alpha)}\,y^{\alpha-1}e^{-\theta y}\,1_{\mathbb{R}_+}(y)\,dy\\
&= \sum_{x\in B_1\cap\mathbb{N}} \left(\dfrac{1}{x!} \int_{B_2\cap\mathbb{R}_+} \dfrac{\theta^{\alpha}}{\Gamma(\alpha)}\,y^{x+\alpha-1}e^{-(\theta+1) y}\,dy\right)\\
&= \sum_{x\in B_1\cap\mathbb{N}} \biggl(\dfrac{\Gamma(x+\alpha)\,\theta^\alpha}{x!\,\Gamma(\alpha)\,(\theta+1)^{x+\alpha}}\\
&\hspace{5em} \times \int_{B_2\cap\mathbb{R}_+} \dfrac{(\theta+1)^{x+\alpha}}{\Gamma(x+\alpha)}\,y^{x+\alpha-1}e^{-(\theta+1) y}\,dy\biggr).
\end{align*} On reconnaît dans cette dernière intégrale la densité d'une
loi Gamma d'indice \(x+\alpha\) et de paramètre d'échelle \(\theta+1\),
qui correspond exactement à la loi conditionnelle de \(Y\) sachant
\(\{X=x\}\) pour \(x\in\mathbb{N}\). En effet, on a d'une part
\begin{align*}
\mathbb{P}_X(B_1) &= \mathbb{P}_{X,Y}(B_1 \times \mathbb{R}) = \sum_{x\in B_1\cap\mathbb{N}} \dfrac{\theta^\alpha}{\Gamma(\alpha)}\, \dfrac{\Gamma(x+\alpha)}{x!\,(\theta+1)^{x+\alpha}},
\end{align*} ce qui donne bien pour tout \(x\in\mathbb{N}\)~:
\begin{align*}
\mathbb{P}_{Y\mid X=x} (B_2) &= \mathbb{P}\left(Y \in B_2 \mid X = x\right) = \dfrac{\mathbb{P}_{X,Y}\left(\{x\}\times B_2\right)}{\mathbb{P}_X(\{x\})}\\
&= \int_{B_2\cap\mathbb{R}_+} \dfrac{(\theta+1)^{x+\alpha}}{\Gamma(x+\alpha)}\,y^{x+\alpha-1}\,e^{-(\theta+1)y}\,dy.
\end{align*}

\section{Randomisation}\label{III-answer-randomize}

En termes probabilistes et selon les notations de l'exercice, il s'agit
de calculer \(\mathbb{E}(T_N - T_0)\), où la variable aléatoire \(T_N\)
peut s'écrire en fonction d'une somme aléatoire de variables aléatoires
indépendantes~: \[T_N = \sum_{i = 1}^N X_i + T_0.\] Comme la boutique
ferme au bout d'un certain temps, toutes les variables aléatoires
figurant dans l'équation précédente sont bornées, donc intégrables. On
peut ainsi calculer \(\mathbb{E}(T_N-T_0)\) à l'aide de la formule de
l'espérance totale~:
\[\mathbb{E}\left(T_N - T_0\right) = \mathbb{E}\left(\mathbb{E}\left(T_N \mid N\right)\right) - T_0.\]
Pour tout \(n\in\mathbb{N}^\ast\) l'énoncé suggère que \(N\) est
indépendante de \(X_1,\dots,X_n\), elles-mêmes indépendantes et de même
loi que \(X\), d'où~:
\[\mathbb{E}\left(T_n \mid N = n\right) = \sum_{i = 1}^n \mathbb{E}(X_i\mid N = n) = \sum_{i = 1}^n \mathbb{E}(X_i) = n\mathbb{E}(X).\]
Ainsi, en posant
\(\psi : n \in\mathbb{N}^\ast \mapsto n \mathbb{E}(X)\), on obtient
\[\mathbb{E}\left(T_N - T_0\right) = \mathbb{E}\left(\psi(N)\right) - T_0 = \mathbb{E}(N)\mathbb{E}(X) - T_0.\]
C'était prévisible~: en posant arbitrairement \(T_0 = 0\), le temps
d'attente moyen est le temps d'attente moyen entre deux arrivées,
multiplié par le rang moyen du gagnant. Si la loterie dépendait des
temps d'arrivées, par exemple en faisant gagner le premier client qui
arrive au moins 10 minutes après le client précédent, \(\psi\), et donc
le résultat, seraient différents.

\section{Etats cachés --- indépendance
conditionnelle}\label{III-etats-cachuxe9s-induxe9pendance-conditionnelle-1}

\runinhead{Question
1}\phantomsection\addcontentsline{toc}{subsubsection}{Question
1}\label{III-answer-ec1}
Soit \(n\in\mathbb{N}^\ast\). Quels que soient \(x\in\mathbb{R}\) et
\(B\) borélien on a \begin{align*}
\mathbb{P}_{Y_n \mid X_n = x}(B) &= \mathbb{E}\left(1_B(X_n+\epsilon_n) \mid X_n = x \right)\\
&= \int_\mathbb{R}1_B(x+y)\,d\mathbb{P}_{\epsilon_n\mid X_n = x}(y)\\
& = \int_\mathbb{R}1_B(x+y)\,f_\epsilon(y)\,dy \hspace{1em}\text{par indépendance de $X_n$ et $\epsilon_n$}\\
&= \int_B f_\epsilon(y-x)\,dy.
\end{align*} Ainsi, \(\mathbb{P}_{Y_n\mid X_n =x}\) admet bien une
densité~:
\[f_{Y_n\mid X_n=x} : y \in\mathbb{R}\mapsto f_\epsilon(y - x).\]

\runinhead{Question
2}\phantomsection\addcontentsline{toc}{subsubsection}{Question
2}\label{III-answer-ec2}
Soient \(n\in\mathbb{N}^\ast\), \((x_1,\dots,x_n)\in\mathbb{R}^n\) et
\(B_1,\dots,B_n\) des boréliens. Pour simplifier les écritures, on note
\(x_{1:n}\) tout vecteur \((x_1,\dots,x_n)\) de \(\mathbb{R}^n\). Alors
\begin{align*}
&\mathbb{P}_{Y_{1:n}\mid X_{1:n}=x_{1:n}}(B_1\times\dots\times B_n)
=\mathbb{E}\left(\prod_{i = 1}^n 1_{B_i}(X_i + \epsilon_i) \Bigm\vert X_{1:n} = x_{1:n} \right)\\
&= \int_{\mathbb{R}^n} \prod_{i = 1}^n 1_{B_i}(x_i + y_i)\,d\mathbb{P}_{\epsilon_{1:n}\mid X_{1:n} = x_{1:n}}(y_{1:n})\\
&= \int_{\mathbb{R}^n} \prod_{i = 1}^n 1_{B_i}(x_i + y_i)\,d\mathbb{P}_{\epsilon_{1:n}}(y_{1:n}) \ \text{par indépendance des $\epsilon_i$ et $X_j$,}\\
&= \prod_{i = 1}^n \int_{\mathbb{R}} 1_{B_i}(x_i + y_i)\,f_\epsilon(x_i)\,dy_{i} \ \text{par Fubini et indépendance et même loi des $\epsilon_i$,}\\
&= \prod_{i = 1}^n \int_{\mathbb{R}} 1_{B_i}(y_i)\,f_\epsilon(y_i - x_i)\,dy_i\\
& = \prod_{i = 1}^n \int_{\mathbb{R}} 1_{B_i}(y_i)\,f_{Y_i\mid X_i = x_i}(y_i)\,dy_i\ \text{par la question 1,}\\
&= \prod_{i = 1}^n \mathbb{P}_{Y_i\mid X_i = x_i}(B_i).
\end{align*} Les \(n\) relevés de température sont donc bien
indépendants conditionnellement aux états cachés.

\section{Covariance totale}\label{III-answer-covtot}

Tout d'abord, par linéarité de l'espérance conditionnelle on a~:
\begin{align*}
\text{Cov}(X,Y \mid Z) &= \mathbb{E}\Bigl( \bigl( X - \mathbb{E}(X\mid Z) \bigr)\bigl( Y - \mathbb{E}(Y\mid Z) \bigr) \Bigm| Z  \Bigr)\\
&= \mathbb{E}\Bigl( XY - X\mathbb{E}(Y\mid Z) - Y\mathbb{E}(X\mid Z) + \mathbb{E}(X\mid Z)\mathbb{E}(Y\mid Z) \Bigm| Z  \Bigr)\\
&= \mathbb{E}(XY \mid Z) - \mathbb{E}(X\mid Z)\mathbb{E}(Y\mid Z).
\end{align*}

En utilisant la formule de l'espérance totale et la linéarité de
l'espérance, on obtient alors \begin{align*}
\text{Cov}(X,Y) &= \mathbb{E}(XY) - \mathbb{E}(X)\mathbb{E}(Y)\\
&= \mathbb{E}\bigl( \mathbb{E}(XY \mid Z) \bigr) - \mathbb{E}\bigl( \mathbb{E}(X \mid Z) \bigr)\mathbb{E}\bigl( \mathbb{E}(Y \mid Z) \bigr)\\
&= \mathbb{E}\bigl( \mathbb{E}(XY \mid Z) - \mathbb{E}(X\mid Z)\mathbb{E}(Y\mid Z) \bigr)\\
&\ \ \  + \mathbb{E}\bigl( \mathbb{E}(X\mid Z)\mathbb{E}(Y\mid Z) \bigr) - \mathbb{E}\bigl( \mathbb{E}(X \mid Z) \bigr)\mathbb{E}\bigl( \mathbb{E}(Y \mid Z) \bigr)\\
&= \mathbb{E}\bigl(\text{Cov}(X,Y\mid Z)\bigr) + \text{Cov}\bigl( \mathbb{E}(X\mid Z), \mathbb{E}(Y\mid Z) \bigr).
\end{align*}

\section{Débruitage (formule de
Tweedie)}\label{III-duxe9bruitage-formule-de-tweedie-1}

\runinhead{Question
1}\phantomsection\addcontentsline{toc}{subsubsection}{Question
1}\label{III-answer-tweedie1}
\(f_Y(y) = \int_{\mathbb{R}^d} f_{X,Y}(x,y) dx = \int_{\mathbb{R}^d} f_{Y|X=x}(y) f_X(x) dx = \int_{\mathbb{R}^d} f_Z(y-x) f_X(x) dx\)
où
\(f_Z(y-x) = f_{Y|X=x}(y) = (2 \pi)^{-d/2} (\det \Sigma)^{1/2}\exp\left(-\frac12 (y-x)^t\Sigma^{-1}(y-x)\right)\).
On notera que l'on retrouve le produit de convolution de \(f_X\) et
\(f_Z\) comme densité de \(X+Z\).

\runinhead{Question
2}\phantomsection\addcontentsline{toc}{subsubsection}{Question
2}\label{III-answer-tweedie2}
On remarque d'abord que \(\nabla \log f_Y = \frac{\nabla f_Y}{f_Y}\).
Ensuite, par application du résultat de l'exercice B.8.5 du poly de
maths 1 (intégrales à paramètres) dont on vérifie aisément les
hypothèses, on a
\[\nabla f_Y(y) = \int_{\mathbb{R}^d} \nabla f_{Y|X=x}(y) f_X(x) dx\] et
\begin{align*}
\nabla f_{Y|X=x}(y) &= (2 \pi)^{-d/2} (\det \Sigma)^{1/2}\exp\left(-\frac12 (y-x)^t\Sigma^{-1}(y-x)\right) \Sigma^{-1}(x-y)\\
&= f_{Y|X=x}(y) \Sigma^{-1}(x-y)
\end{align*} Finalement,
\[\nabla \log f_Y(y) = \frac{\int_{\mathbb{R}^d} f_{Y|X=x}(y) \Sigma^{-1}(x-y) f_X(x) dx}{f_Y(y)}\]

\runinhead{Question
3}\phantomsection\addcontentsline{toc}{subsubsection}{Question
3}\label{III-answer-tweedie3}
En utilisant l'égalité
\(f_{X,Y}(x,y) = f_{Y|X=x}f_X(x) = f_{X|Y=y}(x)f_Y(y)\), on obtient
\begin{align*}
\nabla \log f_Y(y) &= \frac{\int_{\mathbb{R}^d} f_{Y|X=x}(y) \Sigma^{-1}(x-y) f_X(x) dx}{f_Y(y)}\\
&= \int_{\mathbb{R}^d} f_{X|Y=y}(x) \Sigma^{-1}(x-y) dx \\
& = \Sigma^{-1}\left(\int_{\mathbb{R}^d} x f_{X|Y=y}(x) dx -y \right)
\end{align*} d'où \(\mathbb{E}(X|Y=y) = y + \Sigma\nabla\log f_Y(y)\).

\section{Non-réponse}\label{III-non-ruxe9ponse-1}

\runinhead{Question
1}\phantomsection\addcontentsline{toc}{subsubsection}{Question
1}\label{III-answer-nonrep1}
L'espérance conditionnelle de \(Y\) sachant \(X\) peut s'écrire comme la
solution au problème de minimisation
\[\min_{\phi(X)\in L^2_X} \mathbb{E}\left(\left(Y-\phi(X)\right)^2\right).\]
Or pour \(\phi(X)\in L^2_X\) on a ici
\[\mathbb{E}\left(\left(Y-\phi(X)\right)^2\right) = \mathbb{E}\left( \left(Y - \phi(1)\right)^2 1_{\{1\}}(X) \right) + \mathbb{E}\left( \left(Y - \phi(0)\right)^2 1_{\{0\}}(X) \right),\]
il suffit donc de résoudre pour tout \(x\in\{0,1\}\)
\[\min_{\lambda \in \mathbb{R}} \mathbb{E}\left(\left(Y-\lambda\right)^2 1_{\{x\}}(X)\right).\]
Soit \(x\in\{0,1\}\) et posons
\(J_x : \lambda\in\mathbb{R}\mapsto \mathbb{E}\left(\left(Y-\lambda\right)^2 1_{\{x\}}(X)\right)\).
Alors pour tout \(\lambda\in\mathbb{R}\)
\[J_x(\lambda) = \mathbb{E}\left(Y^21_{\{x\}}(X)\right) + \lambda^2\,\mathbb{P}(X=x) -2\lambda\,\mathbb{E}\left(Y1_{\{x\}}(X)\right)\]
et sa dérivée
\[J_x^\prime(\lambda) = 2\lambda\,\mathbb{P}(X=x) -2\,\mathbb{E}\left(Y1_{\{x\}}(X)\right)\]
s'annule en
\[\lambda_x := \dfrac{\mathbb{E}\left(Y1_{\{x\}}(X)\right)}{\mathbb{P}(X=x)} = \mathbb{E}(Y\mid X = x).\]
On en conclut que
\[\mathbb{E}(Y\mid X) = \mathbb{E}(Y\mid X = 1)1_{\{1\}}(X) + \mathbb{E}(Y\mid X = 0)1_{\{0\}}(X).\]
Or on remarque que \(1_{\{1\}}(X) = X\) et \(1_{\{0\}}(X) = 1 - X\), ce
qui fait de \(\mathbb{E}(Y\mid X)\) une fonction affine de \(X\). Elle
est par définition la meilleure approximation de \(Y\) par une fonction
de \(X\), elle coïncide donc avec l'approximation affine de \(Y\) par
\(X\):
\[\mathbb{E}(Y\mid X) = m + \dfrac{\rho\,\sigma}{\sqrt{p(1-p)}}\,(X - p).\]

\runinhead{Question
2}\phantomsection\addcontentsline{toc}{subsubsection}{Question
2}\label{III-answer-nonrep2}
D'après la question précédente, on a
\(\mathbb{E}(Y\mid X) = m_0 + (m_1 - m_0) X\), la meilleure
approximation affine de \(Y\) par \(X\). Ainsi, \(m_0\) et \(m_1\)
satisfont
\[\left|\begin{array}{l} m_1 - m_0 = \dfrac{\rho\sigma}{\sqrt{p(1-p)}},\\[1em] m_0 = m - (m_1-m_0)p,  \end{array}\right. \Leftrightarrow \left|\begin{array}{l} m_1 = m + \rho\sigma\sqrt{\dfrac{1-p}{p}},\\[1em] m_0 = m - \rho\sigma\sqrt{\dfrac{p}{1-p}}.  \end{array}\right.\]

\runinhead{Question
3}\phantomsection\addcontentsline{toc}{subsubsection}{Question
3}\label{III-answer-nonrep3}
Par la formule de la variance totale et d'après la question 1, on a
\begin{align*}
\sigma^2 &= \mathbb{V}\left(Y\right) = \mathbb{E}\bigl(\mathbb{V}\left(Y\mid X\right)\bigr) + \mathbb{V}\bigl(\mathbb{E}(Y\mid X)\bigr)\\
&= p\,\sigma^2_1 + (1-p)\,\sigma^2_0 + \dfrac{\rho^2\sigma^2}{p\,(1-p)}\mathbb{V}(X)\\
&= p\,\sigma^2_1 + (1-p)\,\sigma^2_0 + \rho^2\sigma^2.
\end{align*} Cette égalité se simplifie et donne bien
\[\sigma^2 = \dfrac{(1-p)\,\sigma_0^2 + p\,\sigma_1^2}{1-\rho^2}.\]

\runinhead{Question
4}\phantomsection\addcontentsline{toc}{subsubsection}{Question
4}\label{III-answer-nonrep4}
Lorsque \(X\) et \(Y\) sont non corrélées, i.e.~\(\rho = 0\), on obtient
\(m_0 = m_1 = m\) puis \(\sigma^2 = (1-p)\,\sigma_0^2 + p\,\sigma_1^2\).
En d'autres termes, \(\mathbb{E}(Y\mid X) = m\) est une variable
aléatoire constante, et
\(\mathbb{E}\bigl(\mathbb{V}(Y\mid X)\bigr) = \sigma^2\). Dans ce cas,
la non-réponse n'affecte pas l'espérance, mais potentiellement la
variance (la dispersion du temps de travail peut être différente chez
les répondants et les non-répondants). Ces deux propriétés sont encore
vraies en cas d'indépendance entre \(X\) et \(Y\), puisque
l'indépendance implique la non corrélation, mais nous avons de plus
\(\mathbb{V}(Y\mid X) = \sigma^2 = \sigma^2_1 = \sigma^2_0\); la
variable aléatoire \(\mathbb{V}(Y\mid X)\) est elle aussi constante.
Cette fois-ci, la dispersion est la même chez les répondants et les
non-répondants~: la non-réponse n'affecte pas la variance.
